% ═══════════════════════════════════════════════════════════════
%  6. Ablation Studies
% ═══════════════════════════════════════════════════════════════

\section{Ablation Studies}
\label{sec:ablation}

We conducted several targeted experiments to isolate the contribution of individual components and understand where the system's strengths and weaknesses originate.

% ─── 6.1 LoRA vs Zero-Shot ──────────────────────────────────────────────────
\subsection{Effect of LoRA Fine-Tuning on Hallucination}
\label{sec:ablation_lora}

One of the most practical benefits of LoRA fine-tuning turned out to be hallucination suppression. When we ran the base Whisper-Large-v3-Turbo model (without any adapter) on Gujarati agriculture-news audio, it produced fabricated content that did not exist in the input:

\begin{quote}
\textit{Zero-shot Turbo}: ``\ldots EPMC \ldots 500 people swept away by floodwaters \ldots'' \\ \small (Hallucinated---none of this was in the audio.)
\end{quote}

The same audio processed with the LoRA-adapted model produced a concise, grammatical transcription with no ghost text:

\begin{quote}
\textit{Turbo + LoRA(gu)}: A clean Gujarati transcription mentioning ``APMC Market'' and rain damage, with no invented content.
\end{quote}

This suggests that even a small amount of in-domain fine-tuning (160 samples) teaches the model what Gujarati speech ``sounds like'' well enough to suppress the tendency to hallucinate English phrases when the acoustic signal is ambiguous.

% ─── 6.2 In-Domain vs OOD ───────────────────────────────────────────────────
\subsection{In-Domain vs.\ Out-of-Distribution Accuracy}
\label{sec:ablation_ood}

\begin{figure}[t]
\centering
\begin{tikzpicture}
\begin{axis}[
    ybar,
    bar width=18pt,
    width=0.85\columnwidth,
    height=5.5cm,
    ylabel={WER (\%)},
    symbolic x coords={Marathi, Gujarati},
    xtick=data,
    ymin=0, ymax=90,
    legend style={at={(0.5,0.97)}, anchor=north, legend columns=2, font=\small},
    nodes near coords,
    nodes near coords style={font=\scriptsize},
    every node near coord/.append style={anchor=south},
    enlarge x limits=0.4,
    grid=major,
    grid style={gray!20},
]
\addplot[fill=darkgreen!60, draw=darkgreen] coordinates {(Marathi, 29.3) (Gujarati, 29.3)};
\addplot[fill=mbzuaiblue!60, draw=mbzuaiblue] coordinates {(Marathi, 72.1) (Gujarati, 75.0)};
\legend{In-domain (IndicSpeech val), Out-of-distribution (FLEURS)}
\end{axis}
\end{tikzpicture}
\caption{WER comparison between in-domain (IndicSpeech validation split) and out-of-distribution (FLEURS test split) evaluation. The $\sim$43--46 point gap shows that domain mismatch, not model capacity, dominates the absolute error.}
\label{fig:ood_gap}
\end{figure}

Figure~\ref{fig:ood_gap} illustrates the gap between in-domain and out-of-distribution evaluation. On the IndicSpeech validation split, both languages achieve approximately 29.3\% WER---a respectable number given only 160 training samples. On FLEURS, the same models score 72.1\% and 75.0\%. The $\sim$43--46 point difference is entirely attributable to domain shift: different speakers, different vocabulary (news vs.\ dictation), different recording conditions, and different sentence structures.

This finding has a clear practical implication: the most impactful way to improve the system is not architectural changes or hyperparameter tuning, but simply \textbf{training on more diverse data}. Common Voice, YouTube audio, and other natural speech corpora would close this gap substantially.

% ─── 6.3 Memory Efficiency ──────────────────────────────────────────────────
\subsection{Memory Efficiency}
\label{sec:ablation_memory}

\begin{figure}[t]
\centering
\begin{tikzpicture}
\begin{axis}[
    ybar,
    bar width=30pt,
    width=0.7\columnwidth,
    height=5cm,
    ylabel={GPU Memory (GB)},
    symbolic x coords={Separate Models, Shared + LoRA},
    xtick=data,
    ymin=0, ymax=6,
    nodes near coords,
    nodes near coords style={font=\small\bfseries},
    every node near coord/.append style={anchor=south},
    enlarge x limits=0.5,
]
\addplot[fill=red!40, draw=red!70] coordinates {(Separate Models, 4.80)};
\addplot[fill=darkgreen!50, draw=darkgreen] coordinates {(Shared + LoRA, 1.93)};
\end{axis}
\end{tikzpicture}
\caption{GPU memory comparison: three separate Whisper-Turbo instances vs.\ one shared backbone with three LoRA adapters. The adapter approach uses 60\% less memory.}
\label{fig:memory}
\end{figure}

The LoRA architecture provides a significant memory advantage (Figure~\ref{fig:memory}). Deploying three separate Whisper-Turbo models---one per language---would require 4.80\,GB of GPU memory. Our approach loads a single base model plus three lightweight adapters (113\,MB each), totaling 1.93\,GB. This 60\% reduction is not the headline contribution of the work, but it makes the system practical for deployment on consumer-grade GPUs and shared cluster environments.

% ─── 6.4 Training Convergence ────────────────────────────────────────────────
\subsection{Training Convergence}
\label{sec:ablation_training}

\begin{figure}[t]
\centering
\begin{tikzpicture}
\begin{axis}[
    width=0.85\columnwidth,
    height=5.5cm,
    xlabel={Training Step},
    ylabel={Cross-Entropy Loss},
    xmin=0, xmax=550,
    ymin=-0.1, ymax=4.0,
    legend style={at={(0.97,0.97)}, anchor=north east, font=\small},
    grid=major,
    grid style={gray!20},
    thick,
]
\addplot[color=mbzuaiblue, mark=*, mark size=1.5pt] coordinates {
    (10, 3.26) (50, 2.8) (100, 1.5) (200, 0.4) (300, 0.08) (400, 0.01) (500, 0.0018)
};
\addplot[color=red!70!black, mark=square*, mark size=1.5pt] coordinates {
    (10, 3.65) (50, 3.1) (100, 1.8) (200, 0.5) (300, 0.1) (400, 0.015) (500, 0.0016)
};
\legend{Marathi LoRA, Gujarati LoRA}
\end{axis}
\end{tikzpicture}
\caption{Training loss convergence for the two LoRA adapters. Both reach near-zero loss by step 500, indicating clean convergence on the IndicSpeech data.}
\label{fig:training_loss}
\end{figure}

Both LoRA adapters converge smoothly during training (Figure~\ref{fig:training_loss}). Marathi drops from a loss of 3.26 to 0.0018 by step 500, and Gujarati from 3.65 to 0.0016. The clean convergence confirms that 160 training samples is sufficient for the LoRA rank and learning rate we chose. There are no signs of instability or overfitting on this small dataset, likely because LoRA's low-rank constraint acts as an implicit regularizer.

% ─── 6.5 Per-Stage Latency ──────────────────────────────────────────────────
\subsection{Per-Stage Latency Breakdown}
\label{sec:ablation_latency}

When we include TTS in the pipeline, the latency profile shifts dramatically. The autoregressive Parler-TTS model dominates wall-clock time, accounting for roughly 60--70\% of the total speech-to-speech latency. The STT stage (VAD + LID + transcription) takes 1.4--6.9\,s depending on language and audio length, the LLM translation adds roughly 1\,s, and TTS adds another 4--5\,s per utterance.

This breakdown suggests that the single most impactful optimization for interactive use would be \textbf{streaming TTS}---replacing the autoregressive synthesis with a streaming decoder that can start playing audio before the full utterance is generated. Alternatively, a distillation-based non-autoregressive TTS model would eliminate this bottleneck entirely.
